\section{Separating the graph from the two frontiers}\label{four:sec:bounds}

Throughout this section the prepared geometry, its thresholds and collars,
and $\eps\in(0,1)$ are fixed.  We consider arbitrary smooth solutions of
the four stages of the homotopy in Section~\ref{four:sec:system}, on
$\Omega_R$, satisfying $t\ge-\eps$.  The admissible radii satisfy
$R\ge R_{\min}(N)$, where $R_{\min}(N)\to\infty$ and is large enough for
the outer-boundary conclusion of Proposition~\ref{four:prop:floor-exclusion}.
Every assertion of uniformity includes the homotopy parameter and $R$.
As in Remark~\ref{four:rem:constants}, $\Pi_N$ denotes a polynomial bound in
$N$, whereas $C(N)$ can depend arbitrarily on the fixed integer $N$.
We shall keep this distinction through the boundary-flux estimate.

\subsection{The limiting heights and the two total regions}

The large factor $l/\eta_N$ forces a nonconstant test level of $h$ to
approach a hypersurface whose expansion is controlled by the trace
equation.  A sublevel of the relaxed limiting height need not have a
regular boundary.  We first establish the expansion inequality for
every smooth exterior support, including at a plateau and at a face
of the opposite color.

\begin{proposition}[Separation from the two total regions]
\label{four:prop:height-exclusion}
Let $A\subset\overline\Omega$ be compact.  In the first two homotopy
stages the following statements hold.
\begin{enumerate}[label=\textup{(\roman*)}]
\item If $A\cap\cD_{c_b}=\varnothing$, there exist $\xi_A>0$ and
$N_A$ such that $h\ge b_-+\xi_A$ on $A$ for every $N\ge N_A$.
\item If $A\cap\cW_{c_w}=\varnothing$, there exist $\xi_A>0$ and
$N_A$ such that $h\le b_+-\xi_A$ on $A$ for every $N\ge N_A$.
\end{enumerate}
Here the total regions are the full closed regions of
Proposition~\ref{four:prop:geometric-preparation}, including components whose
frontiers do not border $\Omega$.
\end{proposition}

\begin{proof}
We give the lower-height argument first.  Consider any sequence
$N_i\to\infty$ of the indicated solutions; its radii tend to infinity.
Extend $h_i$ by the constant $b_+$ into a small collar across each white
face.  The functions on these enlarged neighborhoods are continuous.
Define their lower relaxed limit by
\[
 \underline h(x)=\liminf_{\substack{i\to\infty\\y\to x}}h_i(y).
\]
It is lower semicontinuous and locally bounded by
Lemma~\ref{four:lem:coarse-height}.  Fix levels
\begin{equation}\label{four:bnd:levels}
 b_-<k<k'<\min\{0,b_-+\xi_1\},\qquad k'+a_0<0,
\end{equation}
where $\xi_1$ is the height-cutoff constant in $D_0$.
We initially work away from black faces.

Suppose a smooth $\phi$ touches $\underline h$ from below at $x$,
with $\underline h(x)<k'$ and $\dd\phi(x)\ne0$.  Subtracting a small
fourth-order function of distance makes the contact strict on a
small closed neighborhood without changing its second jet.  The
definition of the relaxed limit and minimization on this neighborhood
give, after passing to a subsequence, points $x_i\to x$ and constants $c_i$ such that $\phi+c_i$
touches $h_i$ from below at $x_i$, and
$h_i(x_i)\to\underline h(x)$.  These contacts are interior contacts of
the original exterior: their heights are eventually strictly below
$k'<b_+$, so they cannot occur on a white face or in its constant
extension.  This remains true when $x$ itself is on a white face.

At contact, $\dd h_i=\dd\phi$ and
$\Hess h_i\ge\Hess\phi$.  The positive trace matrix therefore gives
\begin{equation}\label{four:bnd:test-trace}
 \tr_{A_{\chi_i}}K+
 \frac{1}{\sqrt{(\eta_{N_i}/l_i)^2+|\dd\phi|^2}}
       \tr_{A_{\chi_i}}\Hess\phi
 \le h_i+C+\alpha_iD_0,
 \qquad 0\le\alpha_i\le1.
\end{equation}
Here $\alpha_i=0$ in the first stage.  The cutoff construction gives
$D_0\le0$ at these low heights, everywhere: its white contribution
vanishes there, and its black contribution is nonpositive.  Moreover
\[
 \frac{l_i}{\eta_{N_i}}\ge
 \frac{\ell_i}{\eta_{N_i}}=e^{\eps N_i/2}\longrightarrow\infty.
\]
Consequently $|w_i|\to1$, $d_i\to0$, $\chi_i\to0$,
and $A_{\chi_i}$ tends to the orthogonal projection onto
$(\grad\phi)^\perp$.  Passing to the
limit in \eqref{four:bnd:test-trace} proves
\begin{equation}\label{four:bnd:limiting-expansion}
 \Theta_K[\phi]:=
 \tr_{(\grad\phi)^\perp}K+
 \frac{\tr_{(\grad\phi)^\perp}\Hess\phi}{|\dd\phi|}
 \le k'+a_0.
\end{equation}
The normal here points toward increasing $\phi$.

We next transfer this test inequality to the closed set
$D_k=\{\underline h\le k\}$ by
Lemma~\ref{n3:apriori:sublevel-support}, with ambient dimension four,
$u=\underline h$, tensor $Q=K$, and $q_0=k'+a_0$.
The coarse height bound makes $u$ finite, its defining relaxed limit
makes it lower semicontinuous, and
Equation~\eqref{four:bnd:limiting-expansion} gives the required bound
for every nonzero-gradient lower test at a height below $k'$.
Away from black faces these hypotheses hold on open neighborhoods;
at white faces use the constant extensions already constructed.
The lemma therefore bounds every exterior support of $D_k$ there
by $k'+a_0$, including when the sublevel has a plateau.

This argument also rules out a boundary layer at a white face.  The
constant extension need not make $\underline h$ continuous there.
Nevertheless all approximating low contacts used above are interior,
so the support inequality holds at a limiting white-face contact.
At every point strictly on the added side the relaxed limit equals
$b_+>k$, so $D_k$ lies locally in the original closed exterior.
The reversed white face would therefore be an exterior support of
$D_k$ with expansion
\[
 -H+P_B=-c_w>0,
\]
contradicting $k'+a_0<0$.  Hence $D_k$ misses every white face.

Adjoin the full black total region and consider
$D_k\cup\cD_{c_b}$ in the filled black barrier manifold.  At a
frontier point belonging to $\partial\cD_{c_b}$, an exterior support
of this union also supports the smooth region $\cD_{c_b}$, so smooth
tangency bounds its expansion by $c_b$.  At other frontier points
the bound is $k'+a_0>c_b$.  The coarse end-height decay makes the
additional sublevel compact.  It lies strictly inside the distant
barrier sphere: otherwise a sphere at its greatest radius is an
exterior support with positive expansion, a contradiction.  There
are no remaining contacts with the boundary of the filled barrier
manifold.  Lemma~\ref{four:lem:support-enclosure} therefore gives
\begin{equation}\label{four:bnd:sublevel-enclosure}
 D_k\cup\cD_{c_b}\subset\cD_{c'}
 \quad\hbox{whenever } k'+a_0<c'<0
\end{equation}
in the chosen threshold range.

If $A$ is disjoint from $\cD_{c_b}$, its distance from that closed
region is positive.  Right continuity first permits $c'>c_b$ close
enough that $A\cap\cD_{c'}=\varnothing$.  We then choose $k,k'$ as
in \eqref{four:bnd:levels}, close enough to $b_-$ that
$k'+a_0<c'$.  It follows that $A\cap D_k=\varnothing$.  If a uniform
eventual lower bound $h\ge k$ on $A$ failed, a sequence of violating
points in the compact set $A$ would give a point of $A\cap D_k$ in
its lower relaxed limit.  This contradiction proves the first claim,
with, for example, $\xi_A=k-b_-$.

For the second claim apply the same proof to
$(-h,-K,-C)$.  The low-height modification is now $-D_0$ and again
has the required nonpositive sign.  Extend across black faces by
the high constant $-b_-$.  At a limiting contact, the reversed black
face has expansion $-c_b>0$ for $-K$, and is excluded before using
the white barrier manifold.  Adjoining $\cW_{c_w}$ and applying
Lemma~\ref{four:lem:support-enclosure} now places the sublevel in
$\cW_{c'}$.  Right continuity at $c_w$ finishes the proof.  This also
covers an empty chosen white region: the distance-cap formulation
of right continuity makes the nearby regions empty if necessary,
and a nonempty enclosed sublevel is then impossible.
\end{proof}

\subsection{Collar comparisons and boundary slopes}

We shorten the disjoint collars to $0\le s\le s_0$ so that the offset
is exactly linear there.  Their outer leaves are a positive distance
from the corresponding closed total region.  Proposition~\ref{four:prop:height-exclusion}
therefore provides a strict height gap on those leaves in the first
two stages.  All these leaves and the compact sets needed below are
fixed before increasing $N$.

\begin{lemma}[Collar slopes]\label{four:bnd:collar-slopes}
For all sufficiently large $N$, there are constants $0<c_4<C_4$,
independent of $N,R$ and the homotopy parameter, such that in the
first two stages
\begin{equation}\label{four:bnd:signed-slopes}
 \frac{c_4}{\eta_N}\le\varepsilon_B\partial_\nu f
 \le\frac{C_4}{\eta_N},\qquad
 \varepsilon_B=\begin{cases}1&\text{on black faces},\\-1&\text{on white faces}.
 \end{cases}
\end{equation}
On these collars $h\ge b_-$ on black faces and their exterior
collars, and $h\le b_+$ on white faces and their exterior collars.
In stage three,
\begin{equation}\label{four:bnd:third-slopes}
 |\partial_\nu f|\le C_4/\eta_N.
\end{equation}
In stage four $f=0$.
\end{lemma}

\begin{proof}
Write $\nu_s=\grad s/|\dd s|$ and let $P_s,H_s$ be the tangential trace
of $K$ and the mean curvature of the collar leaf, with normal $\nu_s$.
For a test height $h_*=b_*+\psi(s)$ of positive slope, comparison
with a solution uses the test gradient and the solution's value of
$Z$.  With the resulting implicit $t$, direct contraction gives
\begin{equation}\label{four:bnd:collar-trace}
 \tr_{\Achi}(K+H^{f_*})
 =P_s+aH_s+\chi K(\nu_s,\nu_s)
 +a\chi\left(
 \frac{\Hess s(\nu_s,\nu_s)}{|\dd s|}
       +|\dd s|\frac{\psi''}{\psi'}\right).
\end{equation}
Indeed $f_*=h_*/\eta_N$ and
$\Hess f_*=(\psi'\Hess s+\psi''\dd s^2)/\eta_N$.
The factor $l\psi'/(\eta_N\sqrt D)$ equals $a/|\dd s|$,
and $\tr_{\nu_s^\perp}\Hess s=|\dd s|H_s$.  For a test
$b_*-\psi(s)$ with $\psi'>0$, all terms with prefactor $a$ in
\eqref{four:bnd:collar-trace} change sign.

If the slope is between two fixed positive constants, the floor
and the fixed collar geometry imply
\begin{equation}\label{four:bnd:collar-smallness}
 d\le C(\eta_N/\ell)^2=C e^{-\eps N},\qquad
 1-a=\frac{d}{1+a}\le d,\qquad
 N\chi=\frac{3N}{3+pv}d\ge d,
\end{equation}
and $a\ge1/2$ for large $N$.
Choose a smooth nonnegative cutoff $\vartheta_2$, equal to one on
$[0,1]$ and zero on $[2,\infty)$, and set
\[
 b_N(s)=A_2N\vartheta_2(Ns),\qquad
 \int_0^{s_0}b_N(s)\,\dd s\le2A_2.
\]
Thus slopes whose logarithmic derivative is $\pm b_N$ remain within
a fixed multiplicative factor $e^{2A_2}$ of their initial slopes.

On a black collar a lower barrier is $b_-+\psi_-(s)$, with
$\psi_-(0)=0$ and $\psi_-''/\psi_-'=b_N$.
First choose $A_2$ large, depending only on bounded collar geometry.
Then choose the slope multiplier sufficiently small that
$\psi_-'\le\kappa_B/2$ throughout the collar and that its value at
$s_0$ is below the solution by the strict outer-leaf gap.
The direction of its gradient is compatible with the face, so
$D_0=0$ at a contact.  Since $P_s+H_s\ge c_b$ and
$C=a_0-\kappa_Bs$, the left side of the trace equation minus its
prescribed right side at this test height is at least
\begin{equation}\label{four:bnd:lower-residual}
 \frac{\kappa_Bs}{2}-Cd+a\chi|\dd s|b_N.
\end{equation}
The constant multiplying $d$ depends on collar geometry, not on
the size of the chosen slope multiplier.  On $s\le1/N$, the last
term dominates $Cd$ by \eqref{four:bnd:collar-smallness} and the choice
of $A_2$.  On $s\ge1/N$, the first term dominates the exponentially
small $Cd$ for sufficiently large $N$.  The residual is therefore
strictly positive.

For an upper barrier use $b_-+\psi_+(s)$ with
$\psi_+''/\psi_+'=-b_N$.  Choose its minimum slope sufficiently
large to dominate the coarse height at $s_0$ and all fixed linear
spatial variations.  Smoothness and $P_0+H_0=c_b$ give
$|P_s+H_s-c_b|\le Cs$.  The residual is consequently bounded above
by
\begin{equation}\label{four:bnd:upper-residual}
 -c s+Cd-a\chi|\dd s|b_N
\end{equation}
for a fixed $c>0$.  It is strictly negative by the same two ranges
of $s$.  In particular, the favorable linear term in these
comparisons comes from the height and the offset; we have not
assumed a positive derivative of the leaf expansion at $s=0$.

At any hypothetical crossing extremum, the test and the solution
have the same gradient and the same $Z$, hence the same $t$ and
positive trace matrix.  The prescribed right side is increasing
in $h$ with derivative at least one.  The strict residual signs,
Hessian comparison, and the endpoint inequalities exclude such
a crossing.  Differentiating the resulting barriers at $s=0$,
and using upper and lower positive bounds for $|\dd s|$, proves
\eqref{four:bnd:signed-slopes} on black faces.  Sign reversal of
$h,K,C$ proves the white assertion.  The compatible-direction
modification again vanishes there.

In stage three let $b_*=(1-\zeta)b$ be the assigned face height.
Use $b_*+\psi(s)$ and $b_*-\psi(s)$, with $\psi'>0$,
$\psi''/\psi'=-b_N$, and with a large fixed minimum slope.
For the limiting coefficients $a=1$, $\chi=0$ at $s=0$ and
height $b_*$, the directional
inequalities in the homotopy construction are
\[
 F_{\rm presc}(\nu)\ge P_B+H,\qquad
 F_{\rm presc}(-\nu)\le P_B-H.
\]
Their residuals have precisely the signs needed for the upper and
lower tests, respectively.  At positive $s$ and finite $d$, smooth
spatial variation and $1-a\le d$, $\chi\le d$ cost at most
$C(s+d)$.  Increasing or decreasing the height by $\psi(s)$
improves the respective residual by at least $\psi(s)$.  Choose
the minimum slope to dominate the $Cs$ error and the coarse height
at the outer collar leaf.  The logarithmic curvature contribution
in \eqref{four:bnd:collar-trace} has the favorable sign for both tests
and dominates $Cd$ on $s\le1/N$; on $s\ge1/N$ the linear height
term dominates it.  The same comparison proves
\eqref{four:bnd:third-slopes}.  All constants are uniform in
$\zeta\in[0,1]$ by smoothness of the prescribed convex combination.
The assertion for stage four is the scalar maximum-principle
conclusion already proved in the homotopy construction.
\end{proof}

Combining the own-color collar comparison with
Proposition~\ref{four:prop:height-exclusion} on the compact complement of
smaller own-color collars yields
\begin{equation}\label{four:eq:height-interval}
 b_-\le h\le b_+
 \quad\text{on }\overline\Omega\cap\supp K
 \quad\text{in the first two stages, for all sufficiently large }N.
\end{equation}
The finitely many compact sets used here are fixed first.  Thus
\eqref{four:eq:height-interval} has exactly the range required in
\eqref{four:eq:trace-slack}; it was not used to establish the floor.

\subsection{Polynomial weighted traces and flux}

The slope sign supplies a trace inequality with a small coefficient
in the boundary flux.  Its proof uses weighted divergences, so it
does not require an upper bound for $Z$ or $|\dd f|$.
All measures in the next lemma are background $g$ measures.

\begin{lemma}[Weighted trace inequalities]\label{four:lem:weighted-trace}
Let $\mathcal C$ be the union of fixed collars, slightly enlarged
if necessary.  In stages one and two, for every nonnegative smooth
$\varphi$ and every sufficiently large $N$,
\begin{align}
 \int_B L_0\varphi\,\dd A_g
 &\le\Pi_N\int_{\mathcal C}u
       \bigl((1+\sqrt{\cT})\varphi+|\dd\varphi|_{\Achi}\bigr)\,\dd V_g,
       \label{four:bnd:weighted-trace}\\
 \int_B\varphi\,\dd A_g
 &\le\Pi_N\int_{\mathcal C}\sqrt D
       \bigl((1+\sqrt{\cT})\varphi+|\dd\varphi|_{\Achi}\bigr)\,\dd V_g.
       \label{four:bnd:unweighted-trace}
\end{align}
On $B$, the corresponding homotopy flux satisfies
\begin{equation}\label{four:eq:boundary-flux-bound}
 |(V_\lambda)_\nu|\le Cud
 \le C\frac{\eta_N}{\ell}L_0.
\end{equation}
The constants $C$ and the coefficients of $\Pi_N$ are independent
of $N,R$ and the solution.
\end{lemma}

\begin{proof}
Set $W=\sqrt d\,w$.  For every covector $\xi$,
\begin{equation}\label{four:bnd:W-dual}
 |\xi(W)|\le|\xi|_{\Ad}
 \le\sqrt{(N+3)/3}\,|\xi|_{\Achi},\qquad
 uW=L_0w,\qquad \sqrt D\,W=w.
\end{equation}
The first assertion follows along the axis from
$|\xi(W)|^2=dv\xi(e)^2\le d\xi(e)^2$ and is immediate
transversely.

For completeness, write $Q_f=\tr_{\Ad}H^f$.  Differentiating $w$
and $D$ gives
\begin{align*}
 \nabla_iw_j&=H^f_{ij}-H^f(w,\partial_i)w_j+p d\,t_iw_j,\\
 \divg w&=Q_f+p d\,w(t),
 &Q_f&=F+(d-\chi)j-\tr_{\Ad}K.
\end{align*}
The two weighted divergences are exactly
\begin{align}
 D^{-1/2}\divg(\sqrt D\,W)
 &=\sqrt d\bigl(Q_f+p d\,w(t)\bigr),\label{four:bnd:first-divergence}\\
 u^{-1}\divg(uW)
 &=\sqrt d\bigl(Q_f+(p+2+p d)w(t)\bigr).
 \label{four:bnd:second-divergence}
\end{align}
Indeed the second formula follows from $uW=L_0w$ and
$\dd\log L_0=(p+2)\dd t$.  Since
$\sqrt d\,|w(t)|\le|\dd t|_{\Ad}$ and
$0\le d-\chi\le d$, coercivity \eqref{four:eq:coercivity} bounds
both right sides by $\Pi_N(1+\sqrt{\cT})$ on the fixed collars.
No factor $l^{-1}$ or $\eta_N^{-1}$ occurs.

On a face $\varepsilon_Bw_\nu=a\ge1/2$ by
\eqref{four:bnd:signed-slopes}.  Integrate the divergence of
$\varepsilon_B\zeta uW\varphi$ separately on each collar, where
$\zeta=1$ at $B$ and $\zeta=0$ near its outer leaf.  The outward
normal of $\Omega_R$ at $B$ is $-\nu$, so the boundary term is
$-\int_B L_0a\varphi$.  Taking absolute values of the remaining
terms and using \eqref{four:bnd:W-dual}--\eqref{four:bnd:second-divergence}
proves \eqref{four:bnd:weighted-trace}.  Replacing $u$ by $\sqrt D$
proves \eqref{four:bnd:unweighted-trace}.  The derivatives of $\zeta$
cost only a fixed constant.

On the boundary the compatible slope makes $D_0=0$, and the
prescribed heights and offsets give exactly
$F-P_B=\varepsilon_BH$.  Moreover $w_\nu=\varepsilon_Ba$ and
$(\Achi K(w,\cdot))_\nu=\varepsilon_Ba\chi K(\nu,\nu)$.
Writing $q_\lambda=1-(1-\lambda)j_0(t)\in[0,1]$, the flux
formula \eqref{four:eq:homotopy-flux} is therefore
\[
 \frac{(V_\lambda)_\nu}{u}
 =q_\lambda\left[
 -\frac{d}{1+a}H-N_Bd
 -(1-\lambda)\varepsilon_Ba\chi K(\nu,\nu)\right].
\]
This is bounded in absolute value by $Cd$ since $\chi\le d$.
Finally $ud=L_0\sqrt d$ and
$\sqrt d\le C\eta_N/\ell$ at $B$ by
\eqref{four:bnd:signed-slopes}.  This proves
\eqref{four:eq:boundary-flux-bound}.
\end{proof}

\subsection{A high-level integral and graph Sobolev inequality}

From this point, constants used to obtain a bounded range may depend
arbitrarily on fixed $N$.  The only large-$N$ absorption needed first
uses the polynomial constants just proved:
\begin{equation}\label{four:bnd:small-boundary-factor}
 \Pi_N\frac{\eta_N}{\ell}
 =\Pi_N e^{-\eps N/2}\longrightarrow0.
\end{equation}

There exists $Z_0(N)>0$, independent of the solution and $R$, such
that
\begin{equation}\label{four:bnd:high-source}
 \Xi\ge\delta_0\cT+\rho+\tfrac12\delta_0\eta_Na\sigma
       \quad\text{on }\{Z>Z_0(N)\},
 \qquad \Xi\ge\delta_0\cT-C(N)\quad\text{everywhere}.
\end{equation}
To see the first assertion, only the support of $m_0$ matters.
There $-\eps\le t\le-\eps+1/N$, whereas
\[
 D=e^{6(Z-t)},\qquad
 \eta_Na\sigma=\frac{\eta_N}{l}
             \left(\sqrt D-\frac1{\sqrt D}\right).
\]
At fixed $N$ this last expression tends uniformly to infinity with
$Z$, and dominates the bounded penalty $\cP$.  The global lower
bound follows from the bounded weights defining $\cP$.

\begin{lemma}[High-level integral]\label{four:bnd:high-level-integral}
In stage one, choose a smooth nondecreasing $k_0$ which is zero on
$(-\infty,Z_0]$ and one on $[Z_0+1,\infty)$.  For all sufficiently
large fixed $N$,
\begin{equation}\label{four:bnd:high-integral}
 \int_{\Omega_R}u k_0(Z)
       \bigl(\cT+\rho+\eta_Na\sigma\bigr)\,\dd V_g\le C(N).
\end{equation}
\end{lemma}

\begin{proof}
The test $k_0(Z)$ vanishes on $S_R$.  Integration by parts gives
\begin{align*}
 \int u k_0\Xi+3\int u k_0'|\dd Z|_{\Achi}^2
 =-\int_B k_0(V_\lambda)_\nu
  -\lambda\int u k_0'
       \langle K(w,\cdot),\dd Z\rangle_{\Achi}.
\end{align*}
The last integral is absorbed into part of the gradient term plus
$C(N)$.  Indeed it is supported in a fixed compact set and in the
strip $Z_0\le Z\le Z_0+1$, where
\begin{equation}\label{four:bnd:strip-weight}
 u=l e^{3Z-t}\le e^{1+3(Z_0+1)+\eps}.
\end{equation}
Thus the residual strip integral has bounded background volume
and a bounded weight, without any preliminary gradient estimate.

By Lemma~\ref{four:lem:weighted-trace}, the absolute boundary contribution
is at most
\[
 \Pi_N\frac{\eta_N}{\ell}
 \int_{\mathcal C}u
 \bigl((1+\sqrt{\cT})k_0+k_0'|\dd Z|_{\Achi}\bigr).
\]
The fixed collars have $\min_{\mathcal C}\rho>0$.  Consequently
$1+\sqrt{\cT}\le C(\rho+\delta_0\cT)$ there.
Choose $N$ large using \eqref{four:bnd:small-boundary-factor} to absorb
the nondifferentiated term into the positive source in
\eqref{four:bnd:high-source}.  Young's inequality absorbs the differentiated
term into another part of $\int u k_0'|\dd Z|_{\Achi}^2$, with
a bounded remainder by \eqref{four:bnd:strip-weight}.  This proves
\eqref{four:bnd:high-integral}.
\end{proof}

Let $\Gamma_f$ be the ordinary product graph of $f$ in
$(\Omega_R\times\R,g+\dd z^2)$.  Put $E_f=1+\sigma^2$ and
$G_f=\Id-(\sigma^2/E_f)e\otimes e$, the inverse graph metric
on base covectors.  The following comparisons hold without a
bound for $\sigma$:
\begin{equation}\label{four:bnd:graph-comparison}
 \begin{gathered}
 e^{-1}\sqrt D\,\dd V_g\le\dd V_{\Gamma_f}
       =\sqrt{E_f}\,\dd V_g\le\ell^{-1}\sqrt D\,\dd V_g,\\
 \ell^2\Achi\le G_f\le\frac{e^2(N+3)}3\Achi.
 \end{gathered}
\end{equation}
They follow from $\ell\le l\le e$, hence
$\ell^2\le D/E_f\le e^2$, and
$3d/(N+3)\le\chi\le d$.  Thus graph measure and gradient norms
are comparable to $\sqrt D\,\dd V_g$ and the $\Achi$ norm,
with constants depending only on fixed $N$.

On every fixed compact base set $Q$, \eqref{four:bnd:high-integral} implies
\begin{equation}\label{four:bnd:starting-L2}
 \int_{\Gamma_f\vert_Q}e^{2Z}\,\dd V_{\Gamma_f}\le C(N,Q).
\end{equation}
In fact, for $x=\sqrt D\ge1$, the elementary estimate
$x^{2/3}\le C(N)(1+(\eta_N/l)(x-x^{-1}))$ gives
\[
 D^{1/3}\le C(N)(1+\eta_Na\sigma),\qquad
 e^{2Z}\sqrt D=\frac{u}{l}D^{1/3}.
\]
On $\{Z\ge Z_0+1\}$ use \eqref{four:bnd:high-integral} and the
positive minimum of $\rho$ on $Q$.  On its complement the floor
bounds $D$ and $e^{2Z}$, and the fixed base volume suffices.

\begin{lemma}[Sobolev inequality on the graph]
\label{four:bnd:graph-sobolev}
Fix a compact base patch and a slightly larger compact neighborhood.
For every smooth $\varphi$ supported over the patch, allowed to meet
$B$, a stage-one solution satisfies
\begin{equation}\label{four:bnd:sobolev}
 \left(\int_{\Gamma_f}|\varphi|^4\,\dd V_{\Gamma_f}\right)^{1/2}
 \le C(N)\int_{\Gamma_f}
 \left(|\nabla_{\Gamma_f}\varphi|^2+(1+\cT)\varphi^2\right)
     \,\dd V_{\Gamma_f}.
\end{equation}
The constant depends on the fixed neighborhoods, and is independent
of the graph and the outer radius.
\end{lemma}

\begin{proof}
Extend the compact base neighborhood smoothly to a compact
Riemannian manifold and fix a smooth Euclidean isometric embedding
by Nash's theorem~\cite[Theorem~2, p.~59]{Nash1956}.
Its product with a line embeds the product graph isometrically.
The scalar mean curvature of the product graph is, up to its normal
sign,
\[
 H_{\Gamma_f}=
 \frac{\sqrt D}{l\sqrt{E_f}}
 \left[\tr T-\tr_{e^\perp}K+E_f^{-1}(j-K(e,e))\right].
\]
The prefactor is at most $\ell^{-1}$ and
$E_f^{-1}\le e^2d$.  Coercivity therefore bounds its magnitude
by $C(N)(1+\sqrt{\cT})$.  The additional mean-curvature vector
of the fixed embedding has magnitude at most four times the
bound for its second fundamental form.  Thus the Euclidean
mean-curvature vector has the same bound, independently of the
graph slope.

The four-dimensional Michael--Simon $L^1$ Sobolev inequality
\cite{MichaelSimon1973}, in the compact-support form of
\cite[Chapter~4, Section~5, Theorem~5.7, p.~98]{Simon2014}, applied
to $|\varphi|^3$, gives
\begin{align*}
 \left(\int_{\Gamma_f}|\varphi|^4\right)^{3/4}
 \le C(N)\int_{\Gamma_f}
 \left(|\varphi|^2|\nabla_{\Gamma_f}\varphi|
       +(1+\sqrt{\cT})|\varphi|^3\right)
       +C\int_{\partial\Gamma_f}|\varphi|^3.
\end{align*}
We have retained the boundary term.  Its form follows, for each
individual smooth graph, by applying the boundaryless inequality
to cutoffs tending to one up to the face: the integral of their
normal derivatives tends to the face integral.  Artificial patch
edges do not contribute because $\varphi$ vanishes there.  Along
$B$, $f$ is constant on each face, so the graph face measure is
exactly $\dd A_g$.  Apply \eqref{four:bnd:unweighted-trace} to
$|\varphi|^3$ and use \eqref{four:bnd:graph-comparison}.  This bounds
the face contribution by the same two interior expressions.
If $I=\int_{\Gamma_f}|\varphi|^4$, Cauchy--Schwarz now gives
\[
 I^{3/4}\le C(N) I^{1/2}
 \left(\int_{\Gamma_f}
   (|\nabla_{\Gamma_f}\varphi|^2+(1+\cT)\varphi^2)\right)^{1/2}.
\]
For $I>0$ divide by $I^{1/2}$ and square.  The case $I=0$ is
immediate, proving \eqref{four:bnd:sobolev}.
\end{proof}

\subsection{Iteration and a bounded range for all stages}

\begin{proposition}[Global bounded range]\label{four:prop:bounded-range}
For every sufficiently large fixed $N$ there is a finite $C(N)$
such that every above-floor smooth solution in any of the four
homotopy stages satisfies
\begin{equation}\label{four:eq:bounded-range}
 -\eps\le t\le Z\le C(N),\qquad
 |f|+|\grad f|\le C(N)\quad\text{on }\overline\Omega_R.
\end{equation}
The constant is uniform in the allowed radii and homotopy parameters.
\end{proposition}

\begin{proof}
We first treat stage one, where the compactly supported drift can
prevent a direct maximum argument.  Let $\psi$ be a smooth base
cutoff supported in a fixed compact patch, permitted to meet $B$.
For $k\ge k_*(N)\ge1$, test the divergence equation with
$\psi^2e^{2kZ}/L_0$.  We claim
\begin{equation}\label{four:bnd:exponential-energy}
 \int_{\Gamma_f}e^{2kZ}\psi^2
       (\cT+k|\nabla_{\Gamma_f}Z|^2)
 \le C(N)k\int_{\Gamma_f}e^{2kZ}
       (\psi^2+|\dd\psi|_g^2),
\end{equation}
where $C(N)$ is independent of $k$.

Here are the error estimates establishing the claim.  Let
$\dd\mu_D=\sqrt D\,\dd V_g$, and write
$\kappa=K(w,\cdot)$ in this computation.  Since
$\dd\log L_0=(p+2)\dd t$, the exact tested identity is
\begin{align*}
 &\int e^{2kZ}\psi^2
      (\Xi+6k|\dd Z|_{\Achi}^2)\,\dd\mu_D\\
 &=-\int_B\frac{\psi^2e^{2kZ}}{L_0}(V_\lambda)_\nu\,\dd A_g
   -2k\lambda\int e^{2kZ}\psi^2
                   \langle\kappa,\dd Z\rangle_{\Achi}\,\dd\mu_D\\
 &\quad-2\int e^{2kZ}\psi
             \langle3\dd Z+\lambda\kappa,\dd\psi\rangle_{\Achi}\,\dd\mu_D
   +\int e^{2kZ}\psi^2(p+2)
             \langle3\dd Z+\lambda\kappa,\dd t\rangle_{\Achi}\,\dd\mu_D.
\end{align*}
The global lower bound in \eqref{four:bnd:high-source} supplies
$\delta_0\cT$ and costs only $C(N)\psi^2$ on the right.
Coercivity bounds the final integrand, by Young's inequality, by
\[
 \tfrac18\delta_0\psi^2\cT
       +C(N)\psi^2|\dd Z|_{\Achi}^2+C(N)\psi^2.
\]
Choose $k_*(N)$ sufficiently large to absorb this fixed multiple of
the gradient square.  The drift term carrying $2k$ is bounded by
$k\psi^2|\dd Z|_{\Achi}^2+C(N)k\psi^2$.
The cutoff terms cost another fixed fraction of the $k$-weighted
gradient square plus $C(N)k(\psi^2+|\dd\psi|_g^2)$.

For the boundary term, \eqref{four:eq:boundary-flux-bound} gives
$|(V_\lambda)_\nu|/L_0\le C\sqrt d\le C$.
Apply \eqref{four:bnd:unweighted-trace} to $\psi^2e^{2kZ}$.
The resulting integrand, apart from its common factor $e^{2kZ}$,
is bounded by
\[
 C(N)\left[(1+\sqrt{\cT})\psi^2
      +2|\psi||\dd\psi|_{\Achi}
      +2k\psi^2|\dd Z|_{\Achi}\right].
\]
For any fixed coefficient $b=C(N)$ the estimates
\[
 b\sqrt{\cT}\le\tfrac18\delta_0\cT+2b^2/\delta_0,
 \qquad
 2bk|\dd Z|_{\Achi}
       \le k|\dd Z|_{\Achi}^2+b^2k
\]
show that the boundary cost is absorbed with the claimed remainder
linear in $k$.
Leaving fixed positive fractions of the coercive terms and using
\eqref{four:bnd:graph-comparison} proves
\eqref{four:bnd:exponential-energy}.

Set $U=e^Z$.  Applying \eqref{four:bnd:sobolev} to $\psi U^k$ and
using \eqref{four:bnd:exponential-energy} yields
\begin{equation}\label{four:bnd:iteration-step}
 \left(\int_{\Gamma_f}\psi^4U^{4k}\right)^{1/2}
 \le C(N)k^2\int_{\Gamma_f}U^{2k}
                      (\psi^2+|\dd\psi|_g^2).
\end{equation}
Choose nested base patches $Q_r\Subset Q_s$ in a fixed coordinate
neighborhood, or corresponding half patches at $B$, with
$0<s-r\le1$.  Cutoffs have $|\dd\psi|\le C/(s-r)$.
Iterate \eqref{four:bnd:iteration-step} with
$k_j=2^jk_*$ and geometric gaps.  The factor at the $j$th step is
at most
\[
 \bigl[C(N)k_j^2(1+\mathrm{gap}_j^{-2})\bigr]^{1/(2k_j)}.
\]
Its product is finite, since $\sum_j(j+1)/k_j<\infty$.
Tracking the total gap exponent gives
\begin{equation}\label{four:bnd:first-iteration}
 \sup_{Q_r}U\le C(N)(s-r)^{-2/k_*}
    \left(\int_{\Gamma_f\vert_{Q_s}}U^{2k_*}
                         \,\dd V_{\Gamma_f}\right)^{1/(2k_*)}.
\end{equation}
All comparisons here use the same graph measure of the individual
solution; its constants are uniform by the preceding lemmas.

The starting exponent $2k_*$ can exceed the exponent in
\eqref{four:bnd:starting-L2}.  A second, elementary nested-patch argument
closes this point.  Write $M(s)=\sup_{Q_s}U$.  Interpolation gives
\[
 \left(\int_{\Gamma_f\vert_{Q_s}}U^{2k_*}\right)^{1/(2k_*)}
 \le M(s)^{1-1/k_*}
       \left(\int_{\Gamma_f\vert_{Q_s}}U^2\right)^{1/(2k_*)}.
\]
The last integral is uniformly bounded by \eqref{four:bnd:starting-L2}.
If $k_*=1$ this already suffices.  Otherwise Young's inequality
applied to \eqref{four:bnd:first-iteration} gives, for any $q\in(0,1)$,
\begin{equation}\label{four:bnd:second-iteration}
 M(r)\le qM(s)+C(N,q)(s-r)^{-2}.
\end{equation}
Choose increasing radii $r_j$ tending to a fixed larger radius,
with $r_{j+1}-r_j$ proportional to $2^{-j}$, and take $q<1/4$.
Iteration sums a geometric series with ratio $4q<1$.
The remaining $q^jM(r_j)$ tends to zero: for each individual
smooth solution all these patches lie in one compact neighborhood
with finite supremum.  Thus $M(r_0)\le C(N)$ uniformly.  A finite
cover gives a uniform $Z$ bound on a compact set containing $B$
and the support of $K$.

Outside this set $K=0$.  At an interior maximum with $Z>Z_0(N)$,
$\dd Z=0$ and
\[
 \divg(3u\Achi\grad Z)=3u\tr_{\Achi}\Hess Z\le0,
\]
whereas $u\Xi>0$ by \eqref{four:bnd:high-source}.  The outer value
is $Z=0$, so this maximum argument extends the compact bound
through the end.  This proves the stage-one $Z$ estimate.

In stages two and three $\lambda=0$, so the same maximum argument
works everywhere in the interior.  On $B$ the already established
face-gradient bounds show that $Z$ can be high only if $t\ge1$:
if $t<1$, then
\[
 Z=t+\tfrac16\log(1+l^2|\dd f|^2)
 \le1+\tfrac16\log(1+e^2C_4^2\eta_N^{-2}).
\]
For $t\ge1$ the boundary switch satisfies $j_0=1$ and the
prescribed conormal flux in \eqref{four:eq:homotopy-flux} is zero.
A high boundary maximum would contradict the boundary point
principle applied locally on the high-level set to
$\divg(3u\Achi\grad Z)=u\Xi>0$.
For this application the individual smooth solution gives a
smooth positive elliptic matrix; no uniform ellipticity constant
is being assumed before the bound.  Equivalently, the outward
derivative at a nonconstant boundary maximum is strictly positive,
whereas the zero conormal flux forces it to vanish.  This bounds
$Z$ in these two stages.

In stage four $f=0$, hence $t=Z$.  At a positive scale multiplying
the source and boundary data, the same high-level interior and
boundary arguments apply.  At zero scale the homogeneous mixed
problem gives $Z=0$.  The resulting bound is uniform in the scale.

Finally $Z-t=\frac16\log D\ge0$, so the floor and the upper bound
give $D\le e^{6(C(N)+\eps)}$.  Since $l\ge\ell>0$,
$|\dd f|\le\ell^{-1}\sqrt{D-1}\le C(N)$.
The coarse bound for $h=\eta_Nf$ gives $|f|\le C(N)$.
This proves all of \eqref{four:eq:bounded-range}.
\end{proof}

The bounded range now permits fixed-$N$ regularity and exhaustion.
The final mass estimate will use only the earlier polynomial bounds
in Lemma~\ref{four:lem:weighted-trace} and
\eqref{four:eq:boundary-flux-bound}; the graph Sobolev and Moser constants
are not needed in that limit.
